\documentclass[preprint,review,12pt]{elsarticle}

\usepackage{amsmath,amssymb,amsfonts,amsthm}
\usepackage{graphicx}
\usepackage{booktabs}
\usepackage{multirow}
\usepackage{hyperref}
\usepackage{lineno}

\newtheorem{proposition}{Proposition}
\newtheorem{corollary}{Corollary}

\modulolinenumbers[5]

\journal{Results in Engineering}

\begin{document}

\begin{frontmatter}

\title{Bridging the MV/LV Gap: Virtual Slack Nodes and Positional Encodings for GNN-Based Power Flow on Radial Distribution Networks}

\author[ua]{Dmytro~Voitekh\corref{cor1}}
\ead{d.voitekh@gmail.com}
\cortext[cor1]{Corresponding author}

\author[ua]{Anatolii~Tymoshenko}
\ead{timoshag@i.ua}

\affiliation[ua]{organization={Department of Computer Engineering, Open International University of Human Development ``Ukraine''},
                  addressline={Lvivska St, 23},
                  city={Kyiv},
                  postcode={03115},
                  country={Ukraine}}

\begin{highlights}
\item Virtual slack edges reduce effective graph diameter to two hops, cutting voltage MAE by 34\%
\item Combined model (virtual edges + random walk PE + residual GraphSAGE) achieves 43.8\% MAE reduction
\item LV/MV performance gap reduced from 1.41$\times$ ($p=0.010$) to non-significant 1.21$\times$ ($p=0.182$)
\item Physics-informed voltage smoothness loss provides no improvement---negative result reported
\item Evaluated on 10 SimBench grids $\times$ 2,000 scenarios $\times$ 3 seeds (342 experiments total)
\end{highlights}

\begin{abstract}
\textbf{Introduction.} Graph neural networks (GNNs) offer 10--100$\times$ speedups over Newton--Raphson (NR) solvers for AC power flow approximation, but their accuracy degrades on low-voltage (LV) distribution grids relative to medium-voltage (MV) grids. \textbf{Problem.} A preliminary benchmarking of four GNN architectures on ten SimBench distribution grids revealed a 41\% MAE gap between LV and MV grids ($p=0.010$) attributable to the long graph diameter of radial LV feeders (up to 56 hops), which exceeds the receptive field of a standard 4-layer GNN. \textbf{Goal.} Close the MV/LV accuracy gap without abandoning the computational advantage of GNN surrogates. \textbf{Methodology.} Three complementary techniques are evaluated on ten SimBench grids (four MV, six LV) with 2,000 load scenarios each and three random seeds: (1) virtual slack edges that add bidirectional shortcut connections from the slack bus to every node, reducing the effective graph diameter to two hops (formalized in Proposition~\ref{prop:diameter}); (2) random walk positional encodings with $k{=}16$; and (3) residual GraphSAGE with eight layers. \textbf{Results.} The combined model reduces mean voltage magnitude MAE by 43.8\% (Wilcoxon $p=0.002$) and drops the LV/MV ratio from 1.41$\times$ to 1.21$\times$, rendering the gap statistically non-significant on SimBench ($p=0.182$). Virtual slack edges alone contribute a 34\% MAE reduction. Baseline MAE is almost perfectly rank-correlated with graph diameter (Spearman $\rho=+0.95$, $p<10^{-4}$), providing direct evidence for the receptive-field hypothesis. A physics-informed voltage smoothness loss, tested with $\lambda \in \{10, 100, 1000\}$, yields no measurable benefit; a scale analysis explains why. \textbf{Scientific novelty.} The slack bus is reinterpreted as a physically meaningful virtual node, and rank-correlation evidence links baseline error to topological properties diagnosable before training. \textbf{Practical value.} The combined model remains 2--6$\times$ faster than NR while roughly halving LV-grid error, supporting Monte Carlo hosting-capacity and voltage-quality studies where LV accuracy dominates the computational budget. References 36, tables 5, figures 3.
\end{abstract}

\begin{keyword}
Distribution networks \sep Graph neural networks \sep Positional encodings \sep Power flow \sep Virtual nodes \sep Machine learning
\end{keyword}

\end{frontmatter}

\linenumbers

%% ====================================================================
\section{Introduction}
\label{sec:introduction}

The rapid integration of distributed energy resources into distribution networks necessitates frequent power flow calculations for real-time monitoring, voltage regulation, and hosting capacity assessment \cite{b1}. While the Newton--Raphson (NR) method remains the standard solver, its computational cost becomes prohibitive for large-scale Monte Carlo studies and online applications that require thousands of power flow solutions per minute \cite{b2}.

Graph neural networks have emerged as a promising surrogate for power flow, exploiting the natural graph structure of electrical networks to learn the mapping from load injections to bus voltages \cite{b3,b4}; recent surveys document rapid uptake across power-system applications \cite{b33}. Prior studies have demonstrated 4--100$\times$ speedups over NR while maintaining sub-percent accuracy \cite{b5,b6}. However, the application of GNNs to \emph{distribution} networks---particularly low-voltage (LV) grids---has received less attention despite being the setting where computational savings matter most.

In a preliminary benchmarking study, we compared four GNN architectures on ten SimBench distribution grids and identified a systematic performance gap: LV grids exhibited 41\% higher mean absolute error (MAE) than MV grids ($p=0.010$, Mann--Whitney $U$ test). We attributed this gap to two structural properties of radial LV feeders: (1) large graph diameter (up to 56 hops), which far exceeds the receptive field of a 4-layer GNN, and (2) high bridge fraction, meaning that most edges are topologically critical and cannot be bypassed through alternative paths.

This paper proposes three targeted solutions:
\begin{enumerate}
\item \textbf{Virtual slack edges.} We add bidirectional shortcut edges from the slack bus to every node, reducing the effective graph diameter to at most two hops. This directly addresses the receptive field limitation.
\item \textbf{Random walk positional encodings.} We augment node features with $k$-step random walk landing probabilities, providing each node with a topology-aware fingerprint that distinguishes structurally similar but electrically distant positions.
\item \textbf{Residual GraphSAGE.} We deepen the architecture to eight layers with pre-normalization residual connections, expanding the receptive field while mitigating oversmoothing.
\end{enumerate}

We evaluate each component in isolation and in combination across all ten grids. The combined model achieves a 43.8\% mean MAE reduction ($p=0.002$) and, critically, reduces the LV/MV performance ratio from 1.41$\times$ to 1.21$\times$, rendering it statistically non-significant ($p=0.182$). We also report that a physics-informed voltage smoothness loss provides no improvement, which we discuss as a negative result.

%% ====================================================================
\section{Related work}
\label{sec:related}

\subsection{GNNs for power flow}

GNN-based power flow approximation builds on graph convolutional networks \cite{b3}. Donon \emph{et al.} \cite{b4} introduced a message-passing framework that minimizes Kirchhoff violation for AC power flow, while Lin \emph{et al.} \cite{b5} proposed PowerFlowNet, achieving 4--48$\times$ speedups via a combined message-passing and high-order GCN architecture. Ringsquandl \emph{et al.} \cite{b6} showed that GNNs on power grids benefit from exceptionally deep architectures (up to 13 layers), unlike typical graph benchmarks. Lopez-Garcia and Dom\'{i}nguez-Navarro \cite{b29} introduced a typed graph formulation that distinguishes buses and lines as heterogeneous node types. Beinert \emph{et al.} \cite{b27} forecast node-level power flow at transmission grid nodes with GNN plus embedding multi-task learning. Hu \emph{et al.} \cite{b25} proposed an adaptive graph deep learning scheme (EGAT-PFC) that handles varying topologies during operation, while Ghamizi \emph{et al.} \cite{b10} recently extended GNN power flow to unbalanced three-phase distribution systems with PowerFlowMultiNet.

Physics-informed formulations have also appeared: PINCO \cite{b7} incorporates power balance constraints into GNN training for optimal power flow; Zhang \emph{et al.} \cite{b8} proposed a physics-informed line graph neural network; Gao \emph{et al.} \cite{b30} embed power-flow physics directly into the graph convolution kernel to improve robustness to uncertain injections and topology changes; their follow-up work \cite{b31} generalizes the idea to optimal power flow. Deihim \emph{et al.} \cite{b32} use a GNN to produce warm-start estimates that accelerate interior-point solvers for AC-OPF. For state estimation, Ngo \emph{et al.} \cite{b26} combined graph structure with physics-based regularization, and Madbhavi \emph{et al.} \cite{b28} demonstrated scalability of GNN-based estimators on realistic distribution systems. The PowerGraph benchmark \cite{b9} provides standardized GNN evaluation datasets but focuses on transmission networks. Suri and Mangal \cite{b11} combined GraphSAGE with GRU cells for topology-aware temporal prediction. Okoyomon and Goebel \cite{b12} proposed a generalizability framework for GNN-based power flow on SimBench distribution grids, using positional encodings to improve cross-topology transfer, but did not address the MV/LV accuracy gap.

A parallel line of work uses graph-structured learning for control rather than forward prediction: Chen \emph{et al.} \cite{b34} applied multi-agent graph reinforcement learning with a physical-assistance mechanism to fast voltage regulation in PV-rich distribution networks, and Pei \emph{et al.} \cite{b35} studied topology-robust voltage control via multi-task reinforcement learning. Operational risk assessment under evolving grid conditions has also been approached with GNNs \cite{b36}. Most existing work, however, targets transmission networks or optimal power flow. Distribution-level power flow---where radial topology creates the most acute structural challenges for GNNs---remains comparatively underexplored.

\subsection{Positional encodings for GNNs}

Standard message-passing GNNs are limited by the Weisfeiler--Leman isomorphism test and cannot distinguish certain graph structures \cite{b13}. Positional encodings (PEs) address this by augmenting node features with structural information. Dwivedi \emph{et al.} \cite{b14} introduced Laplacian eigenvector PEs for graph transformers, while Li \emph{et al.} \cite{b15} proposed distance encodings including random walk landing probabilities. These techniques have proven effective in molecular property prediction and combinatorial optimization but have not been applied to power flow.

\subsection{Oversmoothing and deep GNNs}

The oversmoothing phenomenon---where node representations converge to a common value as GNN depth increases---is well documented \cite{b16}. Residual connections \cite{b17}, jumping knowledge (JKNet) \cite{b18}, and DropEdge \cite{b19} have been proposed as mitigations. For power flow on radial grids, where the graph diameter can reach 56 hops, the ability to build deep networks without oversmoothing is essential.

\subsection{Virtual nodes in GNNs}

The virtual node technique, introduced by Gilmer \emph{et al.} \cite{b20} for molecular graphs and later adopted in the Open Graph Benchmark \cite{b21}, adds a global node connected to all others, enabling long-range information propagation in a single message-passing step. While effective in molecular and social network domains, this technique has not been explored for power flow. Our approach adapts it to the power systems domain by using the slack bus as the virtual node, which carries physical meaning as the voltage reference.

%% ====================================================================
\section{Methodology}
\label{sec:methodology}

\subsection{Problem formulation}

We formulate AC power flow as a supervised node regression problem on a graph $\mathcal{G} = (\mathcal{V}, \mathcal{E})$, where $\mathcal{V}$ represents buses and $\mathcal{E}$ represents branches (lines and transformers). Each node $i \in \mathcal{V}$ has input features $\mathbf{x}_i$ comprising active and reactive power injections ($P_i$, $Q_i$) and bus type indicators. The target outputs are voltage magnitude $V_{m,i}$ and voltage angle $V_{a,i}$ for each bus.

Given a dataset of $N$ load scenarios solved by the NR method, we train a GNN $f_\theta: \mathcal{G} \rightarrow \mathbb{R}^{|\mathcal{V}| \times 2}$ to approximate the power flow solution. The base loss is the mean squared error:
\begin{equation}
\mathcal{L}_{\text{MSE}} = \frac{1}{N|\mathcal{V}|} \sum_{n=1}^{N} \sum_{i \in \mathcal{V}} \left\| f_\theta(\mathcal{G}_n)_i - \mathbf{y}_i^{(n)} \right\|_2^2
\label{eq:mse}
\end{equation}
where $\mathbf{y}_i^{(n)} = [V_{m,i}^{(n)}, V_{a,i}^{(n)}]^\top$ is the NR solution for bus $i$ in scenario $n$.

\subsection{Virtual slack edges}

The key insight motivating virtual slack edges is that in a radial distribution feeder, every bus is electrically connected to the substation (slack bus) through a unique path. A 4-layer GNN can only aggregate information from nodes within 4 hops, but buses at the end of long LV feeders may be 56 hops from the slack bus, meaning the voltage reference information never reaches them during message passing.

We address this by adding bidirectional edges between the slack bus $s$ and every other node:
\begin{equation}
\mathcal{E}' = \mathcal{E} \cup \{(s, i) : i \in \mathcal{V} \setminus \{s\}\} \cup \{(i, s) : i \in \mathcal{V} \setminus \{s\}\}
\label{eq:virtual_edges}
\end{equation}

\begin{proposition}[Diameter bound]
\label{prop:diameter}
Let $\mathcal{G} = (\mathcal{V}, \mathcal{E})$ be any connected graph with a distinguished vertex $s \in \mathcal{V}$. The augmented graph $\mathcal{G}' = (\mathcal{V}, \mathcal{E}')$ defined by \eqref{eq:virtual_edges} satisfies $\mathrm{diam}(\mathcal{G}') \leq 2$.
\end{proposition}

\emph{Proof.} For any $u, v \in \mathcal{V}\setminus\{s\}$ with $u \neq v$, the path $u \to s \to v$ has length 2 in $\mathcal{G}'$ by construction. Paths involving $s$ itself have length $\leq 1$. $\square$

\begin{corollary}[One-hop reachability of the slack reference]
\label{cor:onehop}
A single message-passing layer on $\mathcal{G}'$ suffices for information to propagate from the slack bus $s$ to every other node.
\end{corollary}

Proposition~\ref{prop:diameter} is the structural guarantee behind our approach: whereas a $d$-layer GNN on the original radial feeder $\mathcal{G}$ requires $d \geq \mathrm{diam}(\mathcal{G})$ to propagate the slack reference to every bus---which can demand up to 56 layers for LV grids in our benchmark---on $\mathcal{G}'$ a single layer always suffices. The added edges carry zero-impedance features, distinguishing them from physical branches. This approach differs from the generic virtual node technique \cite{b20} in that it reuses an existing physically meaningful node rather than introducing a new artificial one.

The number of additional edges is $2(|\mathcal{V}| - 1)$, which for our grids (15--144 buses) adds at most 286 edges. This is a modest increase relative to the existing edge count and does not significantly impact computation time for small grids, although it does increase the per-layer cost for larger networks.

\subsection{Random walk positional encodings}

To provide topology-aware features, we augment each node with random walk positional encodings (RW-PE) \cite{b15}. For a graph with adjacency matrix $\mathbf{A}$ and degree matrix $\mathbf{D}$, the random walk transition matrix is $\mathbf{T} = \mathbf{D}^{-1}\mathbf{A}$. The RW-PE for node $i$ is defined as:
\begin{equation}
\text{RW-PE}_i = \left[\mathbf{T}_{ii}, \mathbf{T}^2_{ii}, \ldots, \mathbf{T}^k_{ii}\right] \in \mathbb{R}^k
\label{eq:rwpe}
\end{equation}
where $\mathbf{T}^j_{ii}$ is the probability that a random walk starting at node $i$ returns to $i$ after $j$ steps. We use $k=16$ steps, which captures both local structure (short return times for high-degree nodes) and mesoscale topology (longer cycles).

We chose RW-PE over Laplacian eigenvector PE for two reasons: (1) RW-PE does not suffer from eigenvector sign ambiguity, which requires additional sign-invariant processing \cite{b14}, and (2) our experiments showed that Laplacian PE with $k=8$ eigenvectors yielded substantially worse results (MAE $3.02 \times 10^{-3}$~p.u.\ vs.\ $0.80 \times 10^{-3}$~p.u.\ for RW-PE), likely because the low-frequency eigenvectors of radial graphs are highly degenerate.

The RW-PE vector is concatenated with the original node features before the first GNN layer. Since the PE depends only on the graph topology and not on the load scenario, it is computed once per grid and reused across all scenarios.

\subsection{Residual GraphSAGE}

To expand the receptive field beyond four hops while avoiding oversmoothing, we adopt a pre-normalization residual architecture. Each layer $l$ of the Residual GraphSAGE computes:
\begin{equation}
\mathbf{h}_i^{(l+1)} = \mathbf{h}_i^{(l)} + \text{SAGE}\!\left(\text{LN}\!\left(\mathbf{h}_i^{(l)}\right), \left\{\text{LN}\!\left(\mathbf{h}_j^{(l)}\right) : j \in \mathcal{N}(i)\right\}\right)
\label{eq:residual}
\end{equation}
where LN denotes layer normalization and SAGE is the GraphSAGE aggregation (mean aggregator followed by a linear projection) \cite{b22}. The residual connection ensures that even with 8 layers, gradient flow is maintained and node features do not collapse.

We evaluated three anti-oversmoothing strategies in a pilot study (Section~\ref{sec:depth}): residual connections, JKNet \cite{b18}, and DropEdge \cite{b19}. Residual connections at depth 8 provided the best accuracy-stability trade-off (Section~\ref{sec:results_ablation}).

\subsection{Physics-informed loss (attempted)}

We experimented with a voltage smoothness regularization term that penalizes large voltage differences between connected buses:
\begin{equation}
\mathcal{L}_{\text{smooth}} = \frac{1}{|\mathcal{E}|} \sum_{(i,j) \in \mathcal{E}} \left(\hat{V}_{m,i} - \hat{V}_{m,j}\right)^2
\label{eq:smooth}
\end{equation}
combined with a bounds penalty:
\begin{equation}
\mathcal{L}_{\text{bounds}} = \frac{1}{|\mathcal{V}|} \sum_{i \in \mathcal{V}} \max\!\left(0, |\hat{V}_{m,i} - 1.0| - 0.05\right)^2
\label{eq:bounds}
\end{equation}

The total loss was $\mathcal{L} = \mathcal{L}_{\text{MSE}} + \lambda (\mathcal{L}_{\text{smooth}} + \mathcal{L}_{\text{bounds}})$ with $\lambda \in \{10, 100, 1000\}$. None of these configurations improved upon the MSE-only baseline. We hypothesize that the voltage smoothness constraint is already implicitly satisfied when training on Z-score normalized NR solutions, as the MSE loss naturally drives predictions toward physically consistent voltage profiles. The bounds constraint is similarly redundant because distribution grid voltages under normal operating conditions rarely approach their limits. We therefore exclude the physics-informed loss from the combined model.

%% ====================================================================
\section{Experimental setup}
\label{sec:setup}

\subsection{Test grids}

We use ten distribution grids from the SimBench benchmark \cite{b23}, comprising four MV grids (20~kV) and six LV grids (0.4~kV). Table~\ref{tab:grids} summarizes their structural properties. The LV grids exhibit substantially larger diameters (up to 56 hops for LV\_rural2) and higher bridge fractions compared with MV grids, confirming the structural basis for the MV/LV performance gap.

\begin{table}[!t]
\renewcommand{\arraystretch}{1.1}
\caption{Structural properties of SimBench distribution grids.}
\label{tab:grids}
\centering
\begin{tabular}{lccccc}
\toprule
Grid & Buses & Edges & Diam. & $R/X$ & Bridge \\
 & & & (hops) & ratio & frac. \\
\midrule
MV\_rural & 97 & 96 & 22 & 1.02 & 0.96 \\
MV\_semiurb & 116 & 119 & 14 & 0.98 & 0.89 \\
MV\_urban & 144 & 147 & 12 & 0.85 & 0.88 \\
MV\_comm & 89 & 88 & 20 & 1.05 & 0.97 \\
LV\_rural1 & 15 & 14 & 13 & 5.21 & 1.00 \\
LV\_rural2 & 96 & 95 & 56 & 4.87 & 1.00 \\
LV\_rural3 & 60 & 59 & 38 & 4.92 & 1.00 \\
LV\_semiurb4 & 44 & 43 & 28 & 3.45 & 1.00 \\
LV\_semiurb5 & 55 & 54 & 32 & 3.68 & 1.00 \\
LV\_urban6 & 40 & 41 & 18 & 2.95 & 0.95 \\
\bottomrule
\end{tabular}
\end{table}

\subsection{Data generation}

For each grid, we generate 2,000 load scenarios by sampling multiplicative factors uniformly from $[0.7, 1.3]$ applied to the nominal active and reactive power loads, representing $\pm 30\%$ load variation. Each scenario is solved using the pandapower NR solver \cite{b24} with standard convergence settings. We apply four critical corrections to the SimBench models: transformer shift degrees set to zero, continuous bus reindexing, NaN voltage-dependent load parameters replaced with zero, and all switches removed. These corrections are essential for NR convergence on the full set of grids.

The data are split into training (70\%), validation (15\%), and test (15\%) sets. Input features (power injections) and targets (voltage magnitude and angle) are Z-score normalized using training set statistics computed per grid.

\subsection{Training protocol}

All models use the AdamW optimizer with learning rate $10^{-3}$, weight decay $10^{-4}$, and OneCycleLR scheduling over 200 epochs. Early stopping with patience 25 on validation loss prevents overfitting. The hidden dimension is 64 for all architectures. We report results averaged over three random seeds (42, 123, 456) with standard deviations.

All experiments run on CPU (Apple M1 Pro) to ensure fair timing comparisons between GNN inference and the NR solver. We measure wall-clock time for both the GNN forward pass and NR convergence on the same hardware.

%% ====================================================================
\section{Results}
\label{sec:results}

\subsection{Ablation study}
\label{sec:results_ablation}

Table~\ref{tab:ablation} summarizes the individual contributions of each proposed component relative to the 4-layer GraphSAGE baseline.

\begin{table}[!t]
\renewcommand{\arraystretch}{1.1}
\caption{Ablation study: individual components vs.\ baseline.}
\label{tab:ablation}
\centering
\begin{tabular}{lccc}
\toprule
Method & MAE $V_m$ & LV/MV & Improv. \\
 & ($\times 10^{-3}$~p.u.) & gap & (\%) \\
\midrule
Baseline (GraphSAGE $d{=}4$) & 0.88 $\pm$ 0.25 & 1.41$\times$ & --- \\
\midrule
E1: Physics loss ($\lambda{=}10$) & 0.88 $\pm$ 0.25 & 1.41$\times$ & ${\sim}$0 \\
E1: Physics loss ($\lambda{=}100$) & 0.88 $\pm$ 0.25 & 1.41$\times$ & ${\sim}$0 \\
E1: Physics loss ($\lambda{=}1000$) & 0.88 $\pm$ 0.25 & 1.41$\times$ & ${\sim}$0 \\
\midrule
E2: RW-PE ($k{=}16$) & 0.80 $\pm$ 0.21 & 1.38$\times$ & 9 \\
E2: Dist-from-slack & 0.81 $\pm$ 0.22 & 1.42$\times$ & 8 \\
E2: Laplacian PE ($k{=}8$) & 3.02 $\pm$ 1.89 & 1.41$\times$ & $-$244 \\
\midrule
E4: Virtual slack edges & 0.58 & 1.20$\times$ & 34 \\
\bottomrule
\end{tabular}
\end{table}

\textbf{E1: Physics-informed loss.} Across all three values of $\lambda$, the physics-informed loss produced no measurable improvement. Averaged over ten grids and three seeds, the mean MAE change relative to baseline was $-0.004 \times 10^{-3}$, $-0.006 \times 10^{-3}$, and $-0.004 \times 10^{-3}$~p.u.\ for $\lambda \in \{10, 100, 1000\}$ respectively, with the largest single-grid deviation only $0.03 \times 10^{-3}$~p.u.---well within seed-to-seed noise and two orders of magnitude smaller than the $0.44 \times 10^{-3}$~p.u.\ improvement achieved by the combined model. This negative result is consistent across all ten grids and three seeds.

\textbf{E2: Positional encodings.} Random walk PE ($k=16$) achieved a modest 9\% improvement in overall MAE and slightly reduced the LV/MV gap from 1.41$\times$ to 1.38$\times$. Distance-from-slack PE performed comparably (8\% improvement) but did not reduce the gap. Laplacian PE degraded performance severely, likely due to eigenvector sign ambiguity and the near-degenerate spectrum of tree-structured graphs.

\textbf{E4: Virtual slack edges.} This was the most effective single component, reducing overall MAE by 34\% and compressing the LV/MV gap from 1.41$\times$ to 1.20$\times$. The MV average dropped from $0.70 \times 10^{-3}$~p.u.\ to $0.52 \times 10^{-3}$~p.u., while the LV average dropped from $0.99 \times 10^{-3}$~p.u.\ to $0.63 \times 10^{-3}$~p.u. The improvement is disproportionately larger for LV grids, confirming that the virtual edges specifically address the long-diameter bottleneck.

\subsubsection{Depth study}
\label{sec:depth}

We conducted a pilot study on two grids (MV\_rural and LV\_rural2) to evaluate three anti-oversmoothing techniques at depths 4, 8, 16, and 32. Table~\ref{tab:depth} reports the average MAE ($\times 10^{-3}$~p.u.) across both grids.

\begin{table}[!t]
\renewcommand{\arraystretch}{1.1}
\caption{Depth study: MAE $V_m$ ($\times 10^{-3}$~p.u.) by technique and depth.}
\label{tab:depth}
\centering
\begin{tabular}{lcccc}
\toprule
Technique & $d{=}4$ & $d{=}8$ & $d{=}16$ & $d{=}32$ \\
\midrule
Residual & 1.02 & \textbf{0.91} & 0.93 & 4.30 \\
JKNet & 1.03 & 0.93 & 0.95 & 1.16 \\
DropEdge & 1.19 & 1.02 & 1.11 & 4.77 \\
\bottomrule
\end{tabular}
\end{table}

Residual connections at $d=8$ achieved the best MAE (0.91), while JKNet provided the most stable performance across depths (least degradation at $d=32$). DropEdge was consistently inferior. At $d=32$, both residual and DropEdge suffered severe oversmoothing, whereas JKNet maintained reasonable accuracy by selectively combining layer representations. Based on these results, we selected residual connections with $d=8$ for the combined model, as it offers the best peak performance at a practical depth.

\subsection{Combined model}

Table~\ref{tab:combined} presents the main result: the combined model (Residual GraphSAGE $d=8$ + RW-PE $k=16$ + virtual slack edges) compared with the baseline across all ten grids.

\begin{table*}[!t]
\renewcommand{\arraystretch}{1.1}
\caption{Per-grid comparison: baseline GraphSAGE vs.\ combined model.}
\label{tab:combined}
\centering
\begin{tabular}{l cc cc cc}
\toprule
& \multicolumn{2}{c}{MAE $V_m$ ($\times 10^{-3}$~p.u.)} & \multicolumn{2}{c}{MAE $V_a$ (deg)} & \multicolumn{2}{c}{Speedup vs.\ NR} \\
\cmidrule(lr){2-3} \cmidrule(lr){4-5} \cmidrule(lr){6-7}
Grid & Baseline & Combined & Baseline & Combined & Baseline & Combined \\
\midrule
MV\_rural & 0.75 $\pm$ 0.02 & 0.60 $\pm$ 0.04 & 0.096 $\pm$ 0.002 & 0.029 & 25.1$\times$ & 6.1$\times$ \\
MV\_semiurb & 0.72 $\pm$ 0.02 & 0.44 $\pm$ 0.02 & 0.089 $\pm$ 0.002 & 0.026 & 23.5$\times$ & 3.4$\times$ \\
MV\_urban & 0.60 $\pm$ 0.03 & 0.30 $\pm$ 0.01 & 0.062 $\pm$ 0.005 & 0.033 & 25.1$\times$ & 1.5$\times$ \\
MV\_comm & 0.75 $\pm$ 0.04 & 0.42 $\pm$ 0.02 & 0.097 $\pm$ 0.003 & 0.031 & 25.1$\times$ & 2.7$\times$ \\
\midrule
LV\_rural1 & 0.70 $\pm$ 0.05 & 0.34 $\pm$ 0.15 & 0.087 $\pm$ 0.006 & 0.037 & 34.0$\times$ & 2.6$\times$ \\
LV\_rural2 & 1.37 $\pm$ 0.05 & 0.75 $\pm$ 0.02 & 0.117 $\pm$ 0.005 & 0.027 & 26.5$\times$ & 2.2$\times$ \\
LV\_rural3 & 1.11 $\pm$ 0.04 & 0.72 $\pm$ 0.01 & 0.077 $\pm$ 0.001 & 0.024 & 24.0$\times$ & 2.5$\times$ \\
LV\_semiurb4 & 0.95 $\pm$ 0.04 & 0.40 $\pm$ 0.05 & 0.052 $\pm$ 0.002 & 0.014 & 29.2$\times$ & 3.1$\times$ \\
LV\_semiurb5 & 1.15 $\pm$ 0.02 & 0.62 $\pm$ 0.02 & 0.069 $\pm$ 0.000 & 0.022 & 24.3$\times$ & 4.1$\times$ \\
LV\_urban6 & 0.68 $\pm$ 0.01 & 0.34 $\pm$ 0.01 & 0.047 $\pm$ 0.002 & 0.012 & 29.4$\times$ & 2.8$\times$ \\
\midrule
\textbf{MV avg} & \textbf{0.70} & \textbf{0.44} & \textbf{0.086} & \textbf{0.030} & \textbf{24.7$\times$} & \textbf{3.4$\times$} \\
\textbf{LV avg} & \textbf{0.99} & \textbf{0.53} & \textbf{0.075} & \textbf{0.023} & \textbf{27.9$\times$} & \textbf{2.9$\times$} \\
\bottomrule
\end{tabular}
\end{table*}

The combined model reduces MAE $V_m$ on every grid without exception. The largest absolute improvement is on LV\_rural2 ($1.37 \rightarrow 0.75$, a 45\% reduction), which has the largest diameter (56 hops) and was the single worst-performing grid for the baseline. The smallest relative improvement is on MV\_rural (20\%), which already had a moderate baseline error. Voltage angle MAE improves even more dramatically, with reductions of 47--74\% across grids.

A Wilcoxon signed-rank test across all ten grids confirms that the improvement is statistically significant ($p = 0.002$), with a mean improvement of 43.8\% $\pm$ 10.0\%.

\subsection{MV/LV gap analysis}

The primary motivation for this work was to reduce the performance gap between MV and LV grids. Fig.~\ref{fig:gap} illustrates the distribution of MAE $V_m$ for MV and LV grids before and after applying the combined model.

\begin{figure}[!t]
\centerline{\includegraphics[width=\columnwidth]{fig1_mvlv_gap_boxplot.pdf}}
\caption{Distribution of MAE $V_m$ for MV and LV grids. Left: baseline GraphSAGE shows a significant gap (Mann--Whitney $p=0.010$). Right: the combined model reduces the gap to non-significance ($p=0.182$).}
\label{fig:gap}
\end{figure}

For the baseline, the LV/MV ratio is 1.41$\times$ (LV avg: $0.99 \times 10^{-3}$~p.u., MV avg: $0.70 \times 10^{-3}$~p.u.), and the difference is statistically significant at the per-seed level (Mann--Whitney $U$ test on 12 MV vs.\ 18 LV observations, $p = 0.010$). After applying the combined model, the ratio drops to 1.21$\times$ (LV avg: $0.53 \times 10^{-3}$~p.u., MV avg: $0.44 \times 10^{-3}$~p.u.), and the gap becomes statistically non-significant ($p = 0.182$) on the SimBench benchmark used here.

\subsection{Structural correlates of the gap}
\label{sec:correlates}

To move from attribution to evidence, we correlate three structural properties of each grid---diameter, bridge fraction, and $R/X$ ratio---with baseline MAE and with the improvement delivered by the combined model (Table~\ref{tab:correlations}).

\begin{table}[!t]
\renewcommand{\arraystretch}{1.1}
\caption{Spearman rank correlations between grid structure and performance ($n=10$).}
\label{tab:correlations}
\centering
\begin{tabular}{lcc}
\toprule
Property & vs.\ Baseline MAE & vs.\ Combined $\Delta$MAE \\
\midrule
Diameter         & $\rho{=}{+}0.95$, $p{<}10^{-4}$ & $\rho{=}{+}0.66$, $p{=}0.038$ \\
Bridge fraction  & $\rho{=}{+}0.74$, $p{=}0.014$  & $\rho{=}{+}0.82$, $p{=}0.004$ \\
$R/X$ ratio      & $\rho{=}{+}0.50$, $p{=}0.138$  & $\rho{=}{+}0.75$, $p{=}0.013$ \\
\bottomrule
\end{tabular}
\end{table}

Diameter is almost perfectly rank-correlated with baseline error ($\rho = +0.95$), which constitutes direct evidence for the receptive-field hypothesis formulated in Section~\ref{sec:methodology}: grids whose longest path exceeds the receptive field of a 4-layer GNN are the ones on which the baseline suffers. Bridge fraction ($\rho = +0.74$) corroborates this, since high-bridge topologies offer no alternative message-passing paths. The $R/X$ ratio, a purely electrical property, correlates more weakly with baseline error once diameter is accounted for.

The second column links structure to the benefit of our method. The absolute improvement grows with both diameter ($\rho = +0.66$) and bridge fraction ($\rho = +0.82$): the grids where the receptive-field bottleneck is most severe are also the ones that gain the most from the virtual slack edges. This is consistent with the 2-hop diameter guarantee established in Proposition~\ref{prop:diameter}---its effect scales with how far the baseline was from the slack reference to begin with.

\subsection{Speed-accuracy trade-off}

The combined model's improved accuracy comes at the cost of reduced inference speed. The baseline GraphSAGE achieves 24--34$\times$ speedup over NR, while the combined model achieves only 1.5--6.1$\times$ speedup. This reduction stems from two factors: (1) the virtual slack edges roughly double the edge count, increasing per-layer computation; and (2) the deeper architecture (8 vs.\ 4 layers) doubles the number of message-passing steps.

\begin{figure}[!t]
\centerline{\includegraphics[width=\columnwidth]{fig2_pareto_speedup_accuracy.pdf}}
\caption{Speed-accuracy Pareto frontier. The combined model (triangles) achieves lower MAE at the cost of reduced speedup. The baseline (circles) offers higher speedup with moderate accuracy.}
\label{fig:pareto}
\end{figure}

Nevertheless, the combined model remains faster than NR on all grids (minimum 1.5$\times$ on MV\_urban), meaning it still provides a net computational benefit. For applications where accuracy is paramount---such as voltage violation detection or hosting capacity analysis---the combined model offers a favorable trade-off. For applications requiring maximum throughput with acceptable accuracy, the baseline remains preferable.

\begin{figure}[!t]
\centerline{\includegraphics[width=\columnwidth]{fig3_per_grid_comparison.pdf}}
\caption{Per-grid MAE $V_m$ comparison between baseline GraphSAGE and the combined model. Error bars indicate $\pm$1 standard deviation over three seeds.}
\label{fig:pergrid}
\end{figure}

%% ====================================================================
\section{Discussion}
\label{sec:discussion}

\subsection{Why virtual slack edges dominate}

The virtual slack edges produced the largest single improvement (34\% MAE reduction), substantially outperforming positional encodings (9\%) and deeper architectures alone (11\% on the pilot grids). This dominance has a clear structural explanation: in a radial feeder, the slack bus is the voltage reference, and voltage drops accumulate along the unique path from slack to each load bus. By providing a direct 2-hop connection from every bus to the slack, the virtual edges allow the GNN to ``see'' the reference voltage in a single message-passing step, regardless of the electrical distance.

This is more effective than simply deepening the network because a deep GNN still processes information sequentially through intermediate nodes, with potential information loss at each layer. The virtual edges provide a shortcut that bypasses this sequential bottleneck entirely.

\subsection{Why physics-informed loss fails}

Physics-informed approaches have shown promise for optimal power flow \cite{b7,b8}, yet our voltage smoothness loss had no effect. A scale analysis explains the result. In SimBench distribution grids, adjacent buses typically differ in voltage magnitude by $|V_i - V_j| \sim 10^{-3}$~p.u.\ under $\pm30\%$ load variation, so $\mathcal{L}_{\text{smooth}}$ on the ground-truth NR solutions is already of order $10^{-6}$. Because we train on Z-score normalized targets, the MSE term operates in the normalized scale while $\mathcal{L}_{\text{smooth}}$ remains anchored to the physical scale; the $\lambda$ values we swept (10, 100, 1000) span three orders of magnitude and still fail to produce a detectable effect, confirming that the regularization is dominated by, and redundant with, the MSE signal on normalized NR solutions.

The key difference from optimal power flow \cite{b7} is that OPF physics constraints encode feasibility requirements absent from the training data. For power flow approximation, where the training targets already satisfy the physics exactly, additional physics constraints provide no benefit.

\subsection{Limitations}

Our grids have at most 144 buses, which is small relative to real-world feeders that may contain thousands of nodes. The virtual slack edge approach adds $O(|\mathcal{V}|)$ edges, which may become computationally significant for very large grids. Second, the speedup reduction from 25$\times$ to 2--6$\times$ limits the practical advantage for latency-sensitive applications. Third, our evaluation is confined to the SimBench benchmark: all reported statistics---including the structural correlations in Section~\ref{sec:correlates}---describe behavior on this particular family of grids and should not be read as generalizable to real distribution networks without further validation. Fourth, all grids are radial; the 2-hop guarantee of Proposition~\ref{prop:diameter} still holds for meshed topologies, but the receptive-field pressure that makes the method especially useful may be weaker there, and we do not test this empirically. Fifth, we do not consider distributed generation or stochastic renewable profiles, which would add variability beyond the $\pm 30\%$ load perturbation used here. Sixth, we do not compare against recent physics-informed alternatives such as PINCO \cite{b7} or line graph formulations \cite{b8}; such comparisons are left for future work.

\subsection{Practical implications}

For distribution system operators running Monte Carlo studies for hosting capacity assessment or probabilistic voltage analysis, a model that is 2--6$\times$ faster than NR while maintaining sub-$10^{-3}$~p.u.\ voltage magnitude accuracy offers tangible benefits. The reduction of the MV/LV performance gap to statistical non-significance on SimBench is particularly valuable, since LV grids are precisely where the computational burden of repeated power flow is greatest due to their larger number and the need for detailed analysis of voltage-sensitive loads; whether the same holds on utility-scale feeders with different topological statistics is an empirical question that we do not answer here.

%% ====================================================================
\section{Conclusion}
\label{sec:conclusion}

We have proposed and evaluated three techniques to improve GNN-based power flow approximation on radial distribution networks: virtual slack edges, random walk positional encodings, and residual GraphSAGE with eight layers. The combined model achieves a 43.8\% mean reduction in voltage magnitude MAE ($p = 0.002$) relative to a standard 4-layer GraphSAGE baseline, evaluated across ten SimBench distribution grids with three random seeds each.

The most important finding is that the LV/MV performance gap, measured on this benchmark, is reduced from 1.41$\times$ ($p = 0.010$) to 1.21$\times$ ($p = 0.182$), rendering it statistically non-significant. The receptive-field hypothesis is supported by a strong rank correlation between graph diameter and baseline MAE ($\rho = +0.95$, $p < 10^{-4}$), and the absolute improvement also grows with diameter and bridge fraction, consistent with Proposition~\ref{prop:diameter}. Virtual slack edges are the single most impactful component, providing a 34\% MAE reduction. We also report that a physics-informed voltage smoothness loss provides no improvement for power flow approximation, in contrast to its effectiveness for optimal power flow.

Future work should investigate scalability to grids with 1,000 or more buses, where the quadratic cost of virtual edges may necessitate sparse approximations such as hierarchical virtual nodes. Transfer learning across grid topologies and the incorporation of time-varying load and generation profiles are additional promising directions.

%% ====================================================================
\section*{CRediT authorship contribution statement}

\textbf{Dmytro Voitekh:} Conceptualization, Methodology, Software, Validation, Formal analysis, Investigation, Data curation, Writing -- original draft, Visualization.
\textbf{Anatolii Tymoshenko:} Supervision, Writing -- review \& editing.

%% ====================================================================
\section*{Declaration of competing interest}

The authors declare that they have no known competing financial interests or personal relationships that could have appeared to influence the work reported in this paper.

%% ====================================================================
\section*{Data availability}

The SimBench grid models used in this study are publicly available at \url{https://simbench.de/en/}. All source code (data generation, training, evaluation, statistical analysis, and figure generation) and trained model checkpoints are released under the MIT license at \url{https://github.com/dvoitekh/gnn-power-flow-distribution} and archived on Zenodo (DOI to be assigned upon acceptance).

%% ====================================================================
\section*{Declaration of generative AI and AI-assisted technologies in the writing process}

During the preparation of this work, the authors used Claude (Anthropic) for code development assistance. The authors reviewed and edited all content and take full responsibility for the publication.

%% ====================================================================
\section*{Acknowledgments}

The authors thank the SimBench consortium for providing the open-access distribution grid benchmark data.

%% ====================================================================
\bibliographystyle{elsarticle-num}

\begin{thebibliography}{36}

\bibitem{b1} Y. Wang, Q. Chen, T. Hong \emph{et al.}, Review of smart meter data analytics: applications, methodologies, and challenges, IEEE Transactions on Smart Grid 10~(3) (2019) 3125--3148.

\bibitem{b2} F. Capitanescu, Critical review of recent advances and further developments needed in AC optimal power flow, Electric Power Systems Research 136 (2016) 57--68.

\bibitem{b3} T.N. Kipf, M. Welling, Semi-supervised classification with graph convolutional networks, in: Proc. 5th Int. Conf. Learning Representations (ICLR), Toulon, France, 2017.

\bibitem{b4} B. Donon, R. Cl\'{e}ment, B. Donnot \emph{et al.}, Neural networks for power flow: graph neural solver, Electric Power Systems Research 189 (2020) 106547.

\bibitem{b5} N. Lin, S. Orfanoudakis, N. Ordonez Cardenas \emph{et al.}, PowerFlowNet: power flow approximation using message passing graph neural networks, International Journal of Electrical Power \& Energy Systems 160 (2024) 110112.

\bibitem{b6} M. Ringsquandl, H. Sellami, M. Hildebrandt \emph{et al.}, Power to the relational inductive bias: graph neural networks in electrical power grids, in: Proc. 30th ACM Int. Conf. Information \& Knowledge Management (CIKM), 2021, pp.~1538--1547.

\bibitem{b7} A. Varbella, D. Briens, B. Gjorgiev \emph{et al.}, Physics-informed GNN for non-linear constrained optimization: PINCO, a solver for the AC-optimal power flow, arXiv preprint arXiv:2410.04818, 2024.

\bibitem{b8} H.-F. Zhang, X.-L. Lu, X. Ding, X.-M. Zhang, Physics-informed line graph neural network for power flow calculation, Chaos 34~(11) (2024) 113123.

\bibitem{b9} A. Varbella, K. Amara, B. Gjorgiev \emph{et al.}, PowerGraph: a power grid benchmark dataset for graph neural networks, in: Advances in Neural Information Processing Systems 37 (NeurIPS) Datasets and Benchmarks Track, Vancouver, Canada, 2024.

\bibitem{b10} S. Ghamizi, J. Cao, A. Ma, P. Rodriguez, PowerFlowMultiNet: multigraph neural networks for unbalanced three-phase distribution systems, IEEE Transactions on Power Systems (2024). DOI: 10.1109/TPWRS.2024.3465088.

\bibitem{b11} D. Suri, M. Mangal, PowerGNN: a topology-aware graph neural network for electricity grids, arXiv preprint arXiv:2503.22721, 2025.

\bibitem{b12} E. Okoyomon, C. Goebel, A framework for assessing the generalizability of GNN-based AC power flow models, in: Proc. 16th ACM Int. Conf. Future and Sustainable Energy Systems (e-Energy), 2025.

\bibitem{b13} K. Xu, W. Hu, J. Leskovec \emph{et al.}, How powerful are graph neural networks?, in: Proc. 7th Int. Conf. Learning Representations (ICLR), New Orleans, LA, USA, 2019.

\bibitem{b14} V.P. Dwivedi, C.K. Joshi, A.T. Luu \emph{et al.}, Benchmarking graph neural networks, Journal of Machine Learning Research 24~(43) (2023) 1--48.

\bibitem{b15} P. Li, Y. Wang, H. Wang \emph{et al.}, Distance encoding: design provably more powerful neural networks for graph representation learning, in: Advances in Neural Information Processing Systems 33 (NeurIPS), 2020, pp.~4465--4478.

\bibitem{b16} Q. Li, Z. Han, X.-M. Wu, Deeper insights into graph convolutional networks for semi-supervised learning, in: Proc. 32nd AAAI Conf. Artificial Intelligence, New Orleans, LA, USA, 2018, pp.~3538--3545.

\bibitem{b17} K. He, X. Zhang, S. Ren \emph{et al.}, Deep residual learning for image recognition, in: Proc. IEEE Conf. Computer Vision and Pattern Recognition (CVPR), Las Vegas, NV, USA, 2016, pp.~770--778.

\bibitem{b18} K. Xu, C. Li, Y. Tian \emph{et al.}, Representation learning on graphs with jumping knowledge networks, in: Proc. 35th Int. Conf. Machine Learning (ICML), Stockholm, Sweden, 2018, pp.~5453--5462.

\bibitem{b19} Y. Rong, W. Huang, T. Xu \emph{et al.}, DropEdge: towards deep graph convolutional networks on node classification, in: Proc. 8th Int. Conf. Learning Representations (ICLR), 2020.

\bibitem{b20} J. Gilmer, S.S. Schoenholz, P.F. Riley \emph{et al.}, Neural message passing for quantum chemistry, in: Proc. 34th Int. Conf. Machine Learning (ICML), Sydney, Australia, 2017, pp.~1263--1272.

\bibitem{b21} W. Hu, M. Fey, M. Zitnik \emph{et al.}, Open Graph Benchmark: datasets for machine learning on graphs, in: Advances in Neural Information Processing Systems 33 (NeurIPS), 2020, pp.~22118--22133.

\bibitem{b22} W.L. Hamilton, R. Ying, J. Leskovec, Inductive representation learning on large graphs, in: Advances in Neural Information Processing Systems 30 (NeurIPS), Long Beach, CA, USA, 2017, pp.~1025--1035.

\bibitem{b23} S. Meinecke, D. Sarajli\'{c}, S.R. Drauz \emph{et al.}, SimBench---a benchmark dataset of electric power systems to compare innovative solutions based on power flow analysis, Energies 13~(12) (2020) 3290.

\bibitem{b24} L. Thurner, A. Scheidler, F. Sch\"{a}fer \emph{et al.}, pandapower---an open-source Python tool for convenient modeling, analysis, and optimization of electric power systems, IEEE Transactions on Power Systems 33~(6) (2018) 6510--6521.

\bibitem{b25} X. Hu, J. Yang, Y. Gao \emph{et al.}, Adaptive power flow analysis for power system operation based on graph deep learning, International Journal of Electrical Power \& Energy Systems 160 (2024) 110166. DOI: 10.1016/j.ijepes.2024.110166.

\bibitem{b26} Q.-H. Ngo, B. L. H. Nguyen, T. Vu, J. Zhang, Physics-informed graphical neural network for power system state estimation, Applied Energy 358 (2024) 122602. DOI: 10.1016/j.apenergy.2023.122602.

\bibitem{b27} D. Beinert, C. Holzh\"{u}ter, J. M. Thomas, S. Vogt, Power flow forecasts at transmission grid nodes using graph neural networks, Energy and AI 14 (2023) 100262. DOI: 10.1016/j.egyai.2023.100262.

\bibitem{b28} R. Madbhavi, B. Natarajan, B. Srinivasan, Graph neural network-based distribution system state estimators, IEEE Transactions on Industrial Informatics 19 (11) (2023) 11383--11393. DOI: 10.1109/TII.2023.3248082.

\bibitem{b29} T. B. Lopez-Garcia, J. A. Dom\'{i}nguez-Navarro, Power flow analysis via typed graph neural networks, Engineering Applications of Artificial Intelligence 117 (2023) 105567. DOI: 10.1016/j.engappai.2022.105567.

\bibitem{b30} M. Gao, J. Yu, Z. Yang, J. Zhao, Physics embedded graph convolution neural network for power flow calculation considering uncertain injections and topology, IEEE Transactions on Neural Networks and Learning Systems (2024). DOI: 10.1109/TNNLS.2023.3287028.

\bibitem{b31} M. Gao, J. Yu, Z. Yang, A physics-guided graph convolution neural network for optimal power flow, IEEE Transactions on Power Systems 39 (1) (2024) 380--390. DOI: 10.1109/TPWRS.2023.3238377.

\bibitem{b32} A. Deihim, D. Apostolopoulou, E. Alonso, Initial estimate of AC optimal power flow with graph neural networks, Electric Power Systems Research 234 (2024) 110782. DOI: 10.1016/j.epsr.2024.110782.

\bibitem{b33} W. Liao, B. Bak-Jensen, J. R. Pillai, Y. Wang, Y. Wang, A review of graph neural networks and their applications in power systems, Journal of Modern Power Systems and Clean Energy 10 (2) (2022) 345--360. DOI: 10.35833/MPCE.2021.000058.

\bibitem{b34} Y. Chen, Y. Liu, J. Zhao, G. Qiu, H. Yin, Z. Li, Physical-assisted multi-agent graph reinforcement learning enabled fast voltage regulation for PV-rich active distribution network, Applied Energy 351 (2023) 121743. DOI: 10.1016/j.apenergy.2023.121743.

\bibitem{b35} Y. Pei, J. Zhao, Y. Yao, F. Ding, Multi-task reinforcement learning for distribution system voltage control with topology changes, IEEE Transactions on Smart Grid 14 (3) (2023) 2481--2484. DOI: 10.1109/TSG.2022.3233766.

\bibitem{b36} Y. Zhang, P. M. Karve, S. Mahadevan, Graph neural networks for power grid operational risk assessment under evolving unit commitment, Applied Energy 380 (2025) 124793. DOI: 10.1016/j.apenergy.2024.124793.

\end{thebibliography}

\end{document}
